\subsection{集的凸包}

定义 2. --- 给定仿射空间 E 的子集 A，我们称包含 A 的所有凸集之交为 A 的凸包，也就是说（II, 第 9 页, prop. 5），它是包含 A 的最小凸集。

命题 8. --- 对仿射空间 E 的凸子集的任意族 \((\mathbf{A}_{\iota})_{\iota \in \mathbf{I}}\)，\(\bigcup_{\iota \in \mathbf{I}}\mathbf{A}_{\iota}\) 的凸包恰是线性组合 \(\sum_{\iota \in \mathbf{I}}\lambda_{\iota}x_{\iota}\) 组成的集，其中 \(x_{\iota}\in \mathbf{A}_{\iota},\lambda_{\iota}\geqslant 0\) 对所有 \(\iota \in \mathrm{I}\) 成立（除有限多个指标外 \(\lambda_{\iota} = 0\)），且 \(\sum_{\iota \in \mathbf{I}}\lambda_{\iota} = 1\)。

用 C 表示这些线性组合组成的集，显然 C 包含在包含所有 \(A_{\iota}\) 的每个凸集中（II, 第 8 页, prop. 1）；另一方面，\(A_{\iota} \subset C\)（对每个 \(\iota\)）。剩下只需证明 C 是凸的。设 \(x = \sum_{i} \lambda_{i} x_{i}\)、\(y = \sum_{i} \mu_{i} y_{i}\) 是 C 的两点，\(\alpha\) 是满足 \(0 < \alpha < 1\) 的一个数，写 \(\gamma_{\iota} = \alpha \lambda_{\iota} + (1 - \alpha)\mu_{\iota}\)（对每个 \(\iota \in \mathbf{I}\)），并用 J 表示 I 中使 \(\gamma_{\iota} \neq 0\) 的指标组成的（有限）集。我们可以写 \(\alpha x + (1 - \alpha)y = \sum_{\iota \in \mathbf{J}}\gamma_{\iota}z_{\iota}\)，其中

\[
z _ {\iota} = \gamma_ {\iota} ^ {- 1} (\alpha \lambda_ {\iota} x _ {\iota} + (1 - \alpha) \mu_ {\iota} y _ {\iota})
\]

属于 \(\mathbf{A}_{\iota}\)（对所有 \(\iota \in \mathbf{J}\)）；但 \(\sum_{\iota \in \mathbf{J}} \gamma_{\iota} = \alpha \sum_{\iota \in \mathbf{I}} \lambda_{\iota} + (1 - \alpha) \sum_{\iota \in \mathbf{I}} \mu_{\iota} = 1\)，于是我们看到 \(\alpha x + (1 - \alpha)y \in \mathbf{C}\)。命题得证。

推论 1. --- E 的子集 A 的凸包与线性组合 \(\sum_{i} \lambda_{i} x_{i}\) 组成的集恒同，其中 \((x_{i})\) 是 A 中点的任意有限族，数 \(\lambda_{i} > 0\)（对所有 \(i\)）且 \(\sum_{i} \lambda_{i} = 1\)。

凸集 A 生成的仿射线性流形 (A, II, § 9.3) 的维数称为 A 的维数。

设 E 是向量空间。E 中集 A 的均衡包的凸包 C 称为 A 的均衡凸包（或对称凸包）；显然它是包含 A 的最小对称凸集；它也是 \(A \cup (-A)\) 的凸包，因为 A 的均衡包的每个点都属于以 a 与 -a（其中 \(a \in A\)）为端点的线段。集 C 与线性组合 \(\sum_{i} \lambda_{i} x_{i}\)（其中 \(x_{i} \in A\) 且 \(\sum_{i} |\lambda_{i}| \leqslant 1\)）组成的集重合；因为显然

这个点集是凸的且包含 A 与 -A；只需证明它包含在 C 中，为此只需考虑满足 \(\mu = \sum_{i} |\lambda_{i}| > 0\) 的那些线性组合；于是我们可以写 \(\sum_{i} \lambda_{i} x_{i} = \mu \cdot \sum_{i} \alpha_{i} y_{i}\)，其中 \(\alpha_{i} = \lambda_{i} / \mu\)，\(y_{i} = x_{i}\)（若 \(\lambda_{i} \geqslant 0\)）；而 \(\alpha_{i} = -\lambda_{i} / \mu\)，\(y_{i} = -x_{i}\)（若 \(\lambda_{i} < 0\)）；显然 \(\sum_{i} \alpha_{i} = 1\)，我们的断言得证。

推论 2. --- 设 f 是仿射空间 E 到仿射空间 F 中的一个仿射线性映射；对 E 的每个子集 A，\(f(A)\) 的凸包是 A 的凸包在 f 下的像。

对线性映射与均衡凸包有类似的命题。

